\documentclass[../../main.tex]{subfiles}

\begin{document}

\section{Helly 定理的 Alaoglu 推广}\label{section:ch15-alaoglu}

令 $X$ 为赋范线性空间, $B$ 是它的闭单位球, 而 $B^{*}$ 是它的对偶空间 $X^{*}$ 的闭单位球. 假设 $X$ 是可分的. 选取 $B$ 的稠密子集 $\{x_{n}\}$, 且对所有 $\psi,\eta\in B^{*}$ 定义
\begin{align*}
\rho(\psi,\eta)=\sum_{n=1}^{\infty}\frac{1}{2^{n}}\cdot\left|(\psi-\eta)(x_{n})\right|.
\end{align*}
则 $\rho$ 是由 $B^{*}$ 上的弱 $^{*}$ 拓扑诱导的度量(见推论\cref{corollary:ch15-metrizable}). 对一个度量空间, 紧性与序列紧性是相同的. 因此 Helly 定理可重述为: 若 $X$ 是可分的赋范线性空间, 则它的对偶空间 $X^{*}$ 的闭单位球 $B^{*}$ 是弱 $^{*}$ 紧的. 我们现在用 Tychonoff 乘积定理证明可分性的假设是不需要的.

回忆 Tychonoff 乘积定理的特殊情形: 令 $\Lambda$ 为任何集. 考虑 $\Lambda$ 上取值于闭有界区间 $[-1,1]$ 的实值函数族 $\mathcal{F}(\Lambda)$. 考虑 $\mathcal{F}(\Lambda)$ 为具有乘积拓扑的拓扑空间. $\mathcal{F}(\Lambda)$ 上的乘积拓扑在 $f\in\mathcal{F}(\Lambda)$ 的基由形如下式的集合构成:
\begin{align*}
\mathcal{N}_{\varepsilon,\lambda_{1},\cdots,\lambda_{n}}(f)=\left\{f^{\prime}\in\mathcal{F}(\Lambda)\ \middle|\ \left|f^{\prime}(\lambda_{k})-f(\lambda_{k})\right|<\varepsilon,\ 1\leqslant k\leqslant n\right\},
\end{align*}
其中 $\varepsilon>0$ 而 $\lambda_{k}$ 属于 $\Lambda$. Tychonoff 乘积定理蕴涵具有乘积拓扑的 $\mathcal{F}(\Lambda)$ 构成的拓扑空间是紧的. 因此有了由乘积拓扑诱导的拓扑, $\mathcal{F}(\Lambda)$ 的每个闭子集 $\mathcal{F}_{0}(\Lambda)$ 也是紧的.

\begin{theorem}[Alaoglu]\label{theorem:ch15-alaoglu}
令 $X$ 为赋范线性空间. 则它的对偶空间 $X^{*}$ 的闭单位球 $B^{*}$ 关于弱 $^{*}$ 拓扑是紧的.
\end{theorem}

\begin{proof}
分别记 $X$ 与 $X^{*}$ 的闭单位球为 $B$ 与 $B^{*}$. 根据前面的讨论, 从 $B$ 到 $[-1,1]$ 的函数组成的具有乘积拓扑的拓扑空间 $\mathcal{F}(B)$ 是紧的.

定义限制映射 $R:B^{*}\to\mathcal{F}(B)$ 为 $R(\psi)=\psi|_{B}$. 我们宣称: (i) $R(B^{*})$ 是 $\mathcal{F}(B)$ 的闭子集; (ii) 限制映射 $R$ 是从具有弱 $^{*}$ 拓扑的 $B^{*}$ 映上具有乘积拓扑的 $R(B^{*})$ 的拓扑同胚. 暂时假定(i)与(ii)已被证明. 根据前面的讨论, Tychonoff 乘积定理告诉我们 $R(B^{*})$ 是紧的. 因此任何拓扑同胚于 $R(B^{*})$ 的空间是紧的. 特别地, 根据(ii), $B^{*}$ 是弱 $^{*}$ 紧的.

剩下来要证明(i)与(ii). 首先观察到 $R$ 是一对一的, 由于对于 $\psi,\eta\in B^{*}$, $\psi\neq\eta$, 存在某个 $x\in B$ 使得 $\psi(x)\neq\eta(x)$, 因此 $R(\psi)\neq R(\eta)$. 直接比较弱 $^{*}$ 拓扑的基本开集与乘积拓扑的基本开集揭示了 $R$ 是 $B^{*}$ 映上 $R(B^{*})$ 的一个同胚. 剩下来要证明 $R(B^{*})$ 关于乘积拓扑是闭的. 令 $f:B\to[-1,1]$ 为 $R(B^{*})$ 的关于乘积拓扑的闭包点. 为证明 $f\in R(B^{*})$, 仅须证明(见习题\cref{exercise:ch15-01})对所有 $u,v\in B$ 与 $\lambda\in\mathbf{R}$ 且 $u+v$ 与 $\lambda u$ 也属于 $B$, 有
\begin{align}
f(u+v)=f(u)+f(v)\text{ 且 }f(\lambda u)=\lambda f(u).\label{align:ch15-linearity}
\end{align}
然而, 对任何 $\varepsilon>0$, $f$ 的弱 $^{*}$ 邻域 $\mathcal{N}_{\varepsilon,u,v,u+v}(f)$ 包含某个 $R(\psi_{k})$, 由于 $\psi_{k}$ 是线性的, 我们有 $|f(u+v)-f(u)-f(v)|<3\varepsilon$. 因此, \eqref{align:ch15-linearity}的第一个等式成立. 第二个等式的证明类似.
\end{proof}

\begin{corollary}\label{corollary:ch15-compact-K}
令 $X$ 为赋范线性空间. 则存在紧 Hausdorff 空间 $K$, 使得 $X$ 线性同构于赋于最大值范数的 $C(K)$ 的线性子空间.
\end{corollary}

\begin{proof}
令 $K$ 为具有弱 $^{*}$ 拓扑的对偶空间的闭单位球. Alaoglu 定理告诉我们 $K$ 是紧的, 而我们从第 14 章的定理 7 推出它是 Hausdorff 的. 定义 $\Phi:X\to C(K)$ 为 $\Phi(x)=J(x)|_{\kappa}$. 由于自然嵌入 $J:X\to X^{**}$ 是等距, $\Phi$ 也是.
\end{proof}

\begin{corollary}\label{corollary:ch15-extreme-point}
令 $X$ 为赋范线性空间. 则它的对偶空间 $X^{*}$ 的闭单位球 $B^{*}$ 有极值点.
\end{corollary}

\begin{proof}
我们把 $X^{*}$ 考虑为具有弱 $^{*}$ 拓扑的局部凸拓扑空间. 根据 Alaoglu 定理, $B^{*}$ 是凸与紧的. Krein-Milman 引理告诉我们 $B^{*}$ 有一个极值点.
\end{proof}

\medskip\par\noindent\textbf{注.} 
Alaoglu 定理没有告诉我们赋范线性空间的对偶的闭单位球关于弱 $^{*}$ 拓扑是序列紧的. 例如, 对 $X=\ell^{\infty}$, $X^{*}$ 的闭单位球 $B^{*}$ 不是弱 $^{*}$ 序列紧的. 的确, 对每个 $n$ 定义为
\begin{align*}
\text{对所有 }\{x_{k}\}\in\ell^{\infty},\quad\psi_{n}(\{x_{k}\})=x_{n}
\end{align*}
的序列 $\{\psi_{n}\}\subseteq B^{*}$ 没有弱 $^{*}$ 收敛的子序列. Alaoglu 定理是 Helly 定理从紧性而非序列紧性观点的推广. 根据 Helly 定理, 若 $X$ 是可分的, $B^{*}$ 是弱 $^{*}$ 序列紧的, 且即将介绍的推论\cref{corollary:ch15-reflexive-dual}告诉我们若 $X$ 是自反的, $B^{*}$ 也是弱 $^{*}$ 序列紧的.


\begin{exercise}\label{exercise:ch15-01}
对赋范线性空间的闭单位球 $B$, 假定函数 $f:B\to[-1,1]$ 具有性质: 只要 $u,v,u+v$ 与 $\lambda u$ 属于 $B$, 则 $f(u+v)=f(u)+f(v)$ 且 $f(\lambda u)=\lambda f(u)$. 证明 $f$ 是 $X$ 上属于 $X^{*}$ 的闭单位球的线性泛函在 $B$ 上的限制.
\end{exercise}

\begin{exercise}\label{exercise:ch15-02}
令 $X$ 为赋范线性空间而 $K$ 是 $X^{*}$ 的有界凸弱 $^{*}$ 闭子集. 证明 $K$ 有极值点.
\end{exercise}

\begin{exercise}\label{exercise:ch15-03}
证明无穷维赋范线性空间的任何非空弱开集关于范数是无界的.
\end{exercise}

\begin{exercise}\label{exercise:ch15-04}
用 Baire 范畴定理与前一个习题证明无穷维 Banach 空间上的弱拓扑是不可用完备度度量量化的.
\end{exercise}

\begin{exercise}\label{exercise:ch15-05}
是否每个 Banach 空间同构于某个 Banach 空间的对偶?
\end{exercise}

\end{document}